\section{A finite-width upper bound at the limiting critical exponent}
\label{sec:finite-width-limiting-upper}

The scalar comparison yields a finite-width bound with an explicit
subexponential prefactor. In this section the Gaussian scale is critical,
$t=1$, and the required clean output margin is $1/2$.
No large-deviation limit is used.
Write $H=L-1$ and let $A_{N,f}$ denote the event
$f(x)z_w(x)>1/2$ on the entire Boolean cube.

\begin{theorem}[Finite-width universal upper bound]
\label{thm:finite-width-limiting-upper}
There is a universal constant $C<\infty$ such that, for every
$d,N,H\ge1$ and every Boolean function $f$,
\begin{equation}\label{eq:finite-width-limiting-upper}
 \mu_{N,1}(A_{N,f})
 \le \min\left\{1,\exp\left[
       -NV_*+\frac{CN}{H}+CH\log(2N)\right]\right\}.
\end{equation}
Here $V_*\in(1,2)$ is the scalar limiting value defined above.
The constant does not depend on input dimension, width, or depth.
In fact, the same one-sided bound holds for
${s z_w(x_0)>1/2}$ at any fixed input $x_0$ and sign
$s\in\{-1,1\}$. The full global-margin event therefore satisfies
\[
 \mu_{N,1}(A_N)
 \le \min\left\{1,2\exp\left[
       -NV_*+\frac{CN}{H}+CH\log(2N)\right]\right\},
 \qquad A_N=\bigcup_f A_{N,f}.
\]
\end{theorem}

We first record a modulus estimate for the scalar value.
Extend $V$ to zero by $V(0)=V_*$.

\begin{lemma}[A uniform scalar-value modulus]
\label{lem:scalar-value-modulus}
For a universal $C_V$, every $0\le q\le q'\le1$ satisfies
\[
 0\le V(q)-V(q')\le C_V\sqrt{q'-q}.
\]
In addition, $0\le V_*-V(q)\le C_Vq$ for $0\le q\le1$.
\end{lemma}
\begin{proof}
The small-state bridge gives $V_*-V(q)\le17q$ for $q\le1/128$.
For larger $q$, use $0\le V(q)\le V_*<2$ to obtain the second
assertion with, for example, constant $256$.

For $0<q<q'$, we also have
\begin{equation}\label{eq:scalar-relative-modulus}
 V(q)-V(q')\le5(q'/q-1).
\end{equation}
Indeed choose a path from $q'$ arbitrarily close to optimal and of
cost less than two; the preceding uniform scalar upper construction
permits this. If it has at least one step, write $\eta$ for its first
preactivation control. Gaussian entropy duality at parameter
$1/(4q')$ gives
\[
 \int z^2\,d\eta(z)
 \le q'\left[4D(\eta\Vert\mathcal N(0,q'))+2\log2\right]
 <10q'.
\]
Use the same control and all the same subsequent states when starting
at $q$. The first entropy changes by exactly
\[
 \frac12\log(q/q')+
  \frac12(1/q-1/q')\int z^2\,d\eta(z)
 \le5(q'/q-1).
\]
If the selected path has zero steps, its cost is $1/(8q')<2$;
the zero-step cost from $q$ is larger by at most
$2(q'/q-1)$. Taking the approximation error to zero proves
\eqref{eq:scalar-relative-modulus}.

Put $h=q'-q$. If $q\ge\sqrt h$, the relative modulus is at most
$5\sqrt h$. Otherwise $q'\le\sqrt h+h\le2\sqrt h$, and the second
assertion gives
$V(q)-V(q')\le V_*-V(q')\le512\sqrt h$.
The zero-width interval and $q=0$ cases are immediate.
\end{proof}

\begin{lemma}[Occupation with an additive prefix budget]
\label{lem:finite-width-additive-occupation}
Suppose $u_0,\ldots,u_H\in(0,1]$, nonnegative costs $D_j$, and
$e\ge0$ satisfy
\[
 D_j\ge u_j-u_{j-1}+cu_{j-1}^2,\qquad
 \sum_{j=1}^iD_j\le C_0u_i+e,\qquad
 \sum_{j=1}^HD_j\le3.
\]
For constants depending only on $c,C_0$, their minimum
$a=\min_j u_j$ satisfies
\[
 H\le\frac{C_1}{a}+\frac{C_1e}{a^2},\qquad
 a\le\frac{C_2}{H}+C_2\sqrt{e/H}.
\]
\end{lemma}
\begin{proof}
The reciprocal telescoping inequality gives, on any interval whose
pre-states are at least $b$,
\[
 c\,\#\{\text{steps}\}
 \le b^{-2}\sum D_j+1/u_{\rm end}-1/u_{\rm start}.
\]
If $a$ exceeds the fixed threshold $1/256$, the total cost bound
gives a universal length bound. Otherwise partition the path at its
last exits from $[0,2^ka]$, stopping when $2^ka\ge1/256$.
The prefix ending at the global minimum has cost at most $C_0a+e$,
so its length is at most $C/a+Ce/a^2$.
In each subsequent dyadic block, remove its first transition. All
remaining pre-states exceed $b=2^ka$, the endpoint lies in $(b,2b]$,
and the block cost is at most $2C_0b+e$. Its remaining length is
therefore at most $C/b+Ce/b^2$. The final tail above $1/256$ has
universal bounded length by the total cost bound, after again removing
one transition. Summing the geometric series, and restoring the
removed transitions, proves the first estimate. Solving its quadratic
inequality for $a$ gives the second.
\end{proof}

For $q>0$ and $0\le y<1$, define the upper-tail scalar cost
\[
 I_+(q,y)=\inf\left\{
 D(\eta\Vert\mathcal N(0,q)):
       \int\tanh^2z\,d\eta(z)\ge y\right\}.
\]
It has the Chernoff dual formula
\begin{equation}\label{eq:upper-tail-scalar-dual}
 I_+(q,y)=\sup_{\lambda\ge0}
 \left\{\lambda y-
   \log\mathbb E\exp[\lambda\tanh^2(\sqrt q Z)]\right\}.
\end{equation}
For $y$ below the reference mean both sides are zero. Otherwise an
exponential tilt with positive parameter realizes mean $y$; the tilted
mean increases continuously from the reference mean to one.
This proves the formula directly, without an asymptotic Cramer theorem.
Two consequences will be used:
\begin{align}
 I_+(q,y)&\ge V(q)-V(y),
       \label{eq:upper-tail-bellman}\\
 I_+(q,y)&\ge y-q+cq^2\qquad(0<q\le1),
       \label{eq:upper-tail-dissipation}
\end{align}
where $c>0$ is the strict scalar-potential constant.
For the first, any feasible control has next state $q'\ge y$,
so the definition and monotonicity of $V$ give
$V(q)\le D(\eta\Vert\mathcal N(0,q))+V(q')
\le D(\eta\Vert\mathcal N(0,q))+V(y)$.
For the second, entropy duality and
$\log\mathbb E\exp[\tanh^2(\sqrt qZ)]\le q-cq^2$
give the result. These inequalities also hold at $y=0$ with the
extension $V(0)=V_*$.

\begin{proof}[Proof of Theorem~\ref{thm:finite-width-limiting-upper}]
After increasing $C$, the theorem is trivial if $N=1$, $H=1$,
or $H>N$, since $V_*<2$. Thus assume $N\ge2$ and $2\le H\le N$,
and put $K=H-1$.
Fix one input $x_0$ and condition on its entire first hidden vector.
Define
\[
 v_j=\frac1N\sum_{i=1}^N h_i^{j+1}(x_0)^2
       \quad(0\le j\le K).
\]
All these numbers belong to $[0,1]$; the conditional initial value
$v_0$ is arbitrary. If $v_0=0$, every later hidden vector at $x_0$
and its readout are zero, so the asserted bound is immediate.
Otherwise, conditional on all preceding layers, the next
preactivations at $x_0$ are independent Gaussians of variance
$v_{j-1}$. Thus $v_j$ is exactly an average of $N$ independent
$\tanh^2$ Gaussian variables, and, for every $y\ge0$,
\[
 \Pr\{v_j\ge y\mid\text{previous layers}\}
 \le\exp[-N I_+(v_{j-1},y)].
\]
When $v_{j-1}=0$, the vector remains zero and the same statement has
the evident degenerate interpretation.
The right-hand-side rate is nonincreasing in its first argument:
$\tanh^2(\sqrt qZ)$ increases pointwise with $q$, so its nonnegative
exponential moments increase with $q$.

Take the grid mesh $h=N^{-4}$ and $J=N^4$ bins
$[0,h],(h,2h],\ldots,((J-1)h,1]$.
For a path of bin indices, let $u_j$ be its upper endpoint and
$l_j=u_j-h$ its lower endpoint, and put $u_0=v_0$.
By conditioning successively on full preceding layers, the probability
of this bin path, together with $s z_w(x_0)>1/2$ for any fixed sign
$s$, is at most
\[
 \exp[-N\mathcal A(u)],\qquad
 \mathcal A(u)=\sum_{j=1}^K I_+(u_{j-1},l_j)+\frac1{8u_K}.
\]
For bins after the first, the actual event has $v_j>l_j$, which is
contained in $v_j\ge l_j$; for the first bin the lower threshold is zero.
The conditional readout variance is $v_K/N\le u_K/N$, giving the
last term. Monotonicity in the previous variance justifies replacing
its actual value by $u_{j-1}$ at every transition.

We bound the action of every grid path. If $\mathcal A(u)>V_*$,
the claimed lower estimate below is automatic. Otherwise put
$D_j=I_+(u_{j-1},l_j)$ and $\widetilde D_j=D_j+h$.
By the modulus lemma, $V(l_i)-V(u_i)\le C_V/N^2$ for every bin.
Bellman's inequality on any suffix, together with
$V(u_K)\le1/(8u_K)$, gives
\[
 \sum_{j=i+1}^KD_j+\frac1{8u_K}
 \ge V(u_i)-\frac{C_VK}{N^2}.
\]
Consequently every prefix obeys
\[
 \sum_{j=1}^i\widetilde D_j
 \le\mathcal A(u)-V(u_i)+\frac{C_VK}{N^2}+Kh
 \le C_Vu_i+e,\qquad
 e=\frac{C_VK}{N^2}+Kh.
\]
Also $\sum_j\widetilde D_j\le V_*+Kh<3$, and
\eqref{eq:upper-tail-dissipation} yields the exact inequality
\[
 \widetilde D_j\ge u_j-u_{j-1}+cu_{j-1}^2.
\]
Lemma~\ref{lem:finite-width-additive-occupation} therefore supplies
some $j\in\{0,\ldots,K\}$ with
\[
 u_j\le\frac{C_1}{K}+C_1\sqrt{e/K}
       \le\frac{C_2}{K}+\frac{C_2}{N}.
\]
The suffix bound at this layer gives
\[
 \mathcal A(u)\ge V(u_j)-\frac{C_VK}{N^2}
 \ge V_*-\frac{C_3}{K}-\frac{C_3}{N}
                    -\frac{C_VK}{N^2}.
\]
There are $J^K=N^{4K}$ bin paths. Summing their probability bounds,
using $1/K\le2/H$, and absorbing
$C_3+C_VK/N+4K\log N$ into $CH\log(2N)$ proves the claimed
one-sided bound conditionally on every first hidden vector with
$v_0>0$. The proof used no restriction on $v_0$ beyond $v_0\le1$;
the prefix inequality at $i=0$ is automatic. Integrating this uniform
bound over the first layer proves it unconditionally.

For a fixed target, $A_{N,f}$ is contained in
$\{f(x_0)z_w(x_0)>1/2\}$. The full global-margin event is contained
in $\{|z_w(x_0)|>1/2\}$, and the two possible signs cost only a
factor of two. This proves both assertions.
\end{proof}

In particular, if $m=2^d$, $N=m\exp(O(\sqrt{d\log d}))$, and
$H=\exp(O(\sqrt{d\log d}))$ is chosen so that $N/H=o(m)$,
then the correction $CN/H+CH\log(2N)$ is $o(m)$.
This statement alone is an upper bound. In fact, the global-event
version rules out matching lower masses with logarithmic error $o(m)$
uniformly over all Boolean functions: summing such lower bounds over
their $2^m$ disjoint events contradicts the displayed upper bound.
Section~\ref{sec:critical-counting-obstruction} makes this obstruction
quantitative. A posterior comparison must retain the function-specific
finite-width terms concealed by the common leading rate.
